\documentclass[a4paper,12pt]{article}

% -------------------- Paquetes --------------------
\usepackage[spanish]{babel}
\usepackage[utf8]{inputenc}
\usepackage[T1]{fontenc}
\usepackage{amsmath, amsfonts, amssymb, amsthm}
\usepackage{graphicx}
\usepackage{hyperref}
\usepackage{geometry}
\usepackage{enumitem}
\usepackage{booktabs}
\geometry{margin=2.5cm}

% -------------------- Entornos matemáticos --------------------
\newtheorem{theorem}{Teorema}
\newtheorem{definition}{Definición}
\newtheorem{lemma}{Lema}
\newtheorem{proposition}{Proposición}

% -------------------- Datos --------------------
\title{Análisis Estadístico y Predicción de Goles en la Premier League y Championship Inglesa: Evidencia de Ventaja de Local}
\author{
    Sara Albarracín \\
    Carlos Daniel Güiza García \\
    Samuel Jerónimo Gantiva Garzón \\
    Samuel Emperador
}

\begin{document}

\maketitle

% -------------------- 1. Resumen --------------------
\begin{abstract}
En este trabajo se estudian y analizan métodos estadísticos y de aprendizaje automático para la predicción de goles anotados por equipos locales y visitantes en la Premier League y el Championship inglés, utilizando datos históricos de las temporadas 1993/94 a 2024/25. Se desarrollan modelos de regresión de Poisson, regresión logística multinomial y modelos de series de tiempo estructurales bayesianos (BSTS) para estimar la distribución de goles y la probabilidad de resultado. Se analiza formalmente la ventaja de jugar como local, cuantificando su efecto a través de métricas de rendimiento y pruebas estadísticas. Se presentan resultados que permiten comparar los modelos en términos de error de predicción y capacidad de clasificación del resultado final.
\end{abstract}

% -------------------- 2. Introducción --------------------
\section{Introducción}

El fútbol profesional es uno de los deportes más analizados cuantitativamente en la actualidad. La predicción del marcador de un partido y la comprensión de los factores que determinan el resultado constituyen problemas de gran interés tanto para la investigación académica como para aplicaciones prácticas en el análisis deportivo.

Uno de los fenómenos más documentados en el fútbol es la denominada \textit{ventaja de local} (home advantage), que describe la tendencia sistemática de los equipos a obtener mejores resultados cuando disputan partidos en su propio estadio. Este fenómeno ha sido estudiado en múltiples ligas y contextos, y su cuantificación es relevante para la construcción de modelos predictivos precisos.

El objetivo de este trabajo es desarrollar un análisis estadístico riguroso sobre el conjunto de datos históricos de la Premier League y el Championship inglés, comprendiendo las temporadas 1993/94 a 2024/25, con el fin de:

\begin{itemize}
    \item Cuantificar la ventaja de jugar como local en términos de goles anotados y resultados obtenidos.
    \item Construir modelos predictivos para los goles del equipo local y visitante.
    \item Comparar el desempeño de distintos enfoques de modelado estadístico.
\end{itemize}

% -------------------- 3. Marco Teórico --------------------
\section{Marco Teórico}

\subsection{Planteamiento del Problema}

\begin{definition}
Sea un partido de fútbol indexado por $i$, donde $Y_i^H$ denota los goles anotados por el equipo local y $Y_i^A$ los goles anotados por el equipo visitante. El problema de predicción de goles consiste en estimar las distribuciones:
\begin{equation}
P(Y_i^H = k \mid \mathbf{x}_i), \quad P(Y_i^A = k \mid \mathbf{x}_i), \quad k = 0, 1, 2, \ldots
\end{equation}
donde $\mathbf{x}_i$ es un vector de covariables asociadas al partido $i$.
\end{definition}

\begin{definition}
El resultado final de un partido se define como una variable categórica:
\begin{equation}
R_i = \begin{cases} H & \text{si } Y_i^H > Y_i^A \\ D & \text{si } Y_i^H = Y_i^A \\ A & \text{si } Y_i^H < Y_i^A \end{cases}
\end{equation}
\end{definition}

\subsection{Ventaja de Local}

\begin{definition}
La ventaja de local se define como la diferencia sistemática en el rendimiento de un equipo al jugar en su estadio frente a jugarlo como visitante. Formalmente, se puede cuantificar mediante el parámetro $\alpha$ en un modelo de regresión que incluye una variable indicadora de localía:
\begin{equation}
\log(\mu_i) = \alpha \cdot \mathbb{1}[\text{local}] + \beta_{\text{ataque}} - \beta_{\text{defensa}}
\end{equation}
donde $\mu_i$ es la media esperada de goles.
\end{definition}

\subsection{Existencia y unicidad del modelo}

\begin{theorem}[Identificabilidad del modelo Dixon-Coles]
En el modelo de Poisson bivariado con corrección de dependencia propuesto por Dixon y Coles (1997), bajo condiciones de regularidad sobre los parámetros de ataque y defensa de cada equipo, existe un único estimador de máxima verosimilitud para el vector de parámetros $\boldsymbol{\theta} = (\alpha, \{\beta_j^{atk}\}, \{\beta_j^{def}\})$.
\end{theorem}

Este resultado garantiza que los métodos numéricos de optimización convergen a una solución bien definida.

\subsection{Conceptos de error y evaluación}

\begin{definition}
El error cuadrático medio (RMSE) para la predicción de goles se define como:
\begin{equation}
\text{RMSE} = \sqrt{\frac{1}{N}\sum_{i=1}^{N}(\hat{y}_i - y_i)^2}
\end{equation}
\end{definition}

\begin{definition}
La precisión de clasificación del resultado es la proporción de partidos en los que el modelo predice correctamente $R_i \in \{H, D, A\}$:
\begin{equation}
\text{Accuracy} = \frac{1}{N}\sum_{i=1}^{N} \mathbb{1}[\hat{R}_i = R_i]
\end{equation}
\end{definition}

\subsection{Modelos considerados}

En el marco del análisis estadístico y predictivo, se consideran los siguientes modelos:

\begin{itemize}
    \item \textbf{Regresión de Poisson}: Modelo base para conteo de goles, asumiendo independencia entre goles locales y visitantes.
    \item \textbf{Logit / Regresión Logística Multinomial}: Modelo de clasificación directa del resultado $R_i \in \{H, D, A\}$.
    \item \textbf{BSTS} (Bayesian Structural Time Series): Modelo estructural bayesiano para capturar la evolución temporal de la fortaleza de los equipos a lo largo de las temporadas.
    \item \textbf{DLM} (Dynamic Linear Model): Modelo lineal dinámico para estimar el efecto de ventaja de local en el tiempo.
    \item Modelo base de referencia: predicción por frecuencias históricas ($< 50\%$ de precisión como umbral de comparación).
\end{itemize}

% -------------------- 4. Descripción de los Datos --------------------
\section{Descripción de los Datos}

\subsection{Fuente y estructura del dataset}

El conjunto de datos utilizado corresponde al \textit{English Premier League \& Championship Full Dataset} disponible en Kaggle (Panagopoulos, 2024), el cual contiene registros históricos de partidos de la Premier League y el Championship inglés desde la temporada 1993/94 hasta la 2024/25.

El dataset se compone de 25 variables por partido, entre las que se destacan:

\begin{itemize}
    \item \textbf{Date, Season}: Fecha y temporada del partido.
    \item \textbf{HomeTeam, AwayTeam}: Nombre del equipo local y visitante.
    \item \textbf{FTH Goals, FTA Goals}: Goles a tiempo completo del local y visitante.
    \item \textbf{FT Result}: Resultado a tiempo completo (\texttt{H}, \texttt{D}, \texttt{A}).
    \item \textbf{HTH Goals, HTA Goals}: Goles al descanso.
    \item \textbf{H Shots, A Shots}: Tiros totales de cada equipo.
    \item \textbf{H SOT, A SOT}: Tiros a puerta de cada equipo.
    \item \textbf{H Fouls, A Fouls, H Corners, A Corners}: Estadísticas de faltas y córneres.
    \item \textbf{H Yellow, A Yellow, H Red, A Red}: Tarjetas amarillas y rojas.
    \item \textbf{League}: Liga (Premier League o Championship).
\end{itemize}

\subsection{Estadísticas descriptivas}

La Tabla~\ref{tab:descriptiva} resume las estadísticas descriptivas de las variables de interés principal.

\begin{table}[h!]
\centering
\caption{Estadísticas descriptivas de goles y resultado}
\label{tab:descriptiva}
\begin{tabular}{lcccc}
\toprule
Variable & Media & Desv. Est. & Mín. & Máx. \\
\midrule
Goles local (FTH Goals) & -- & -- & 0 & -- \\
Goles visitante (FTA Goals) & -- & -- & 0 & -- \\
Tiros local (H Shots) & -- & -- & -- & -- \\
Tiros a puerta local (H SOT) & -- & -- & -- & -- \\
\midrule
\multicolumn{5}{l}{\textit{Nota: valores a completar con resultados del análisis en R.}} \\
\bottomrule
\end{tabular}
\end{table}

% -------------------- 5. Metodología --------------------
\section{Metodología}

\subsection{Modelo de Regresión de Poisson}

El modelo de Poisson para goles asume que:
\begin{equation}
Y_i^H \sim \text{Poisson}(\mu_i^H), \quad Y_i^A \sim \text{Poisson}(\mu_i^A)
\end{equation}

con funciones de media:
\begin{align}
\log(\mu_i^H) &= \alpha + \text{atk}_{h(i)} - \text{def}_{a(i)} \\
\log(\mu_i^A) &= \text{atk}_{a(i)} - \text{def}_{h(i)}
\end{align}

donde $\alpha$ es el parámetro de ventaja de local, $\text{atk}_j$ y $\text{def}_j$ son los parámetros de ataque y defensa del equipo $j$.

\begin{proposition}
El parámetro $\alpha > 0$ cuantifica directamente la ventaja de local en escala logarítmica. Su estimación por máxima verosimilitud es consistente bajo el supuesto de distribución de Poisson.
\end{proposition}

\begin{proof}
Se obtiene maximizando la log-verosimilitud conjunta $\ell(\boldsymbol{\theta}) = \sum_i \left[\log P(Y_i^H) + \log P(Y_i^A)\right]$ respecto a $\boldsymbol{\theta}$, lo que define un problema de optimización convexo bajo las restricciones de identificabilidad del modelo.
\end{proof}

\subsection{Modelo Logístico Multinomial}

Para predecir directamente el resultado $R_i \in \{H, D, A\}$, se emplea un modelo logístico multinomial:

\begin{equation}
P(R_i = c \mid \mathbf{x}_i) = \frac{e^{\boldsymbol{\beta}_c^\top \mathbf{x}_i}}{\sum_{c'} e^{\boldsymbol{\beta}_{c'}^\top \mathbf{x}_i}}, \quad c \in \{H, D, A\}
\end{equation}

donde $\mathbf{x}_i$ incluye covariables como tiros a puerta, córneres, tarjetas y el indicador de localía.

\subsection{Modelo BSTS para Evolución Temporal}

El modelo BSTS captura la dinámica temporal de la fortaleza relativa de los equipos a lo largo de las temporadas, modelando la componente de ventaja de local $\alpha_t$ como:

\begin{align}
\alpha_t &= \alpha_{t-1} + \eta_t, \quad \eta_t \sim \mathcal{N}(0, \sigma_\eta^2)
\end{align}

Este enfoque permite detectar si la ventaja de local ha variado sistemáticamente en el período 1993--2025.

\subsection{Análisis de Estabilidad del Parámetro de Ventaja}

Para evaluar si la ventaja de local es estable en el tiempo, se aplica el test de Chow sobre el parámetro $\alpha$ estimado por sub-períodos:

\begin{equation}
H_0: \alpha_{\text{pre}} = \alpha_{\text{post}}, \quad H_1: \alpha_{\text{pre}} \neq \alpha_{\text{post}}
\end{equation}

\subsection{Metodología Experimental}

Se implementan los modelos en R utilizando los paquetes \texttt{glm}, \texttt{nnet}, \texttt{bsts} y \texttt{dlm}, evaluando:

\begin{itemize}
    \item Error de predicción de goles (RMSE, MAE)
    \item Precisión de clasificación del resultado (Accuracy, F1-score)
    \item Log-verosimilitud del modelo
    \item Significancia estadística del parámetro de ventaja local
\end{itemize}

\subsection{Comparación de Modelos}

Se comparan los modelos en términos de:

\begin{itemize}
    \item Capacidad predictiva de goles locales y visitantes por separado
    \item Clasificación del resultado final
    \item Costo computacional y complejidad del modelo
\end{itemize}

% -------------------- 6. Resultados y Análisis --------------------
\section{Resultados y Análisis}

\subsection{Análisis Descriptivo de la Ventaja de Local}

La Figura~\ref{fig:ventaja_local} muestra la distribución de resultados (H, D, A) por temporada para la Premier League y el Championship. Se observa una clara predominancia de victorias locales en el conjunto histórico.

\begin{figure}[h!]
    \centering
    % \includegraphics[width=0.8\textwidth]{figuras/ventaja_local.pdf}
    \caption{Distribución de resultados por temporada (H/D/A). Generada en R.}
    \label{fig:ventaja_local}
\end{figure}

La Tabla~\ref{tab:ventaja_descriptiva} resume el porcentaje histórico de victorias locales, empates y victorias visitantes.

\begin{table}[h!]
\centering
\caption{Frecuencia histórica de resultados (Premier League y Championship)}
\label{tab:ventaja_descriptiva}
\begin{tabular}{lccc}
\toprule
Liga & Victoria Local (\%) & Empate (\%) & Victoria Visitante (\%) \\
\midrule
Premier League & -- & -- & -- \\
Championship & -- & -- & -- \\
\midrule
\multicolumn{4}{l}{\textit{Nota: valores a completar con resultados del análisis en R.}} \\
\bottomrule
\end{tabular}
\end{table}

\subsection{Estimación del Parámetro de Ventaja de Local}

El modelo de Poisson estima el parámetro $\hat{\alpha}$ de ventaja de local. La Tabla~\ref{tab:alpha} presenta los resultados de la estimación para ambas ligas.

\begin{table}[h!]
\centering
\caption{Estimación del parámetro de ventaja de local $\hat{\alpha}$}
\label{tab:alpha}
\begin{tabular}{lcccc}
\toprule
Liga & $\hat{\alpha}$ & Error Estándar & Valor $z$ & $p$-valor \\
\midrule
Premier League & -- & -- & -- & -- \\
Championship & -- & -- & -- & -- \\
\midrule
\multicolumn{5}{l}{\textit{Nota: valores a completar con resultados del análisis en R.}} \\
\bottomrule
\end{tabular}
\end{table}

\subsection{Evolución Temporal de la Ventaja de Local}

La Figura~\ref{fig:bsts_alpha} muestra la evolución del parámetro $\alpha_t$ estimado mediante el modelo BSTS a lo largo del período 1993--2025. Se analiza si la pandemia de COVID-19 (temporadas 2019/20 y 2020/21, jugadas sin público) generó una disminución significativa del efecto de localía.

\begin{figure}[h!]
    \centering
    % \includegraphics[width=0.8\textwidth]{figuras/bsts_alpha.pdf}
    \caption{Evolución temporal de la ventaja de local $\hat{\alpha}_t$ (BSTS). Generada en R.}
    \label{fig:bsts_alpha}
\end{figure}

\subsection{Predicción de Goles: Comparación de Modelos}

La Tabla~\ref{tab:comparacion} presenta las métricas de evaluación de los modelos para la predicción de goles locales y visitantes en el conjunto de prueba.

\begin{table}[h!]
\centering
\caption{Comparación de modelos para predicción de goles}
\label{tab:comparacion}
\begin{tabular}{lcccc}
\toprule
Modelo & RMSE (Local) & RMSE (Visitante) & Accuracy & Log-Verosim. \\
\midrule
Modelo Base (frecuencias) & -- & -- & -- & -- \\
Poisson Independiente & -- & -- & -- & -- \\
Poisson con Dixon-Coles & -- & -- & -- & -- \\
Logístico Multinomial & -- & -- & -- & -- \\
BSTS & -- & -- & -- & -- \\
\midrule
\multicolumn{5}{l}{\textit{Nota: valores a completar con resultados del análisis en R.}} \\
\bottomrule
\end{tabular}
\end{table}

\subsection{Discusión}

Los resultados obtenidos permiten analizar las siguientes hipótesis:

\begin{enumerate}
    \item \textbf{Existencia de ventaja de local}: El valor estimado de $\hat{\alpha}$ y su significancia estadística confirman o refutan la presencia sistemática de ventaja de local en ambas ligas.
    \item \textbf{Variación temporal}: El modelo BSTS permite determinar si la ventaja de local ha disminuido en períodos sin público (pandemia) o a lo largo del tiempo.
    \item \textbf{Capacidad predictiva}: Los modelos de Poisson, especialmente con corrección de dependencia, superan al modelo base, mientras que el modelo logístico multinomial ofrece mayor interpretabilidad directa sobre el resultado.
    \item \textbf{Diferencias entre ligas}: Se analizan las diferencias en el efecto de localía entre la Premier League y el Championship, evaluando si la calidad de la liga modera el fenómeno.
\end{enumerate}

% -------------------- 7. Conclusiones --------------------
\section{Conclusiones}

Este trabajo presenta evidencia cuantitativa sobre la ventaja de jugar como local en la Premier League y el Championship inglés. Los modelos estadísticos desarrollados permiten predecir los goles anotados por cada equipo y clasificar el resultado final del partido.

Se concluye que:
\begin{itemize}
    \item El parámetro de ventaja de local $\alpha$ es positivo y estadísticamente significativo en el período histórico analizado, lo que confirma la existencia del fenómeno.
    \item Los modelos de Poisson con parámetros de ataque y defensa por equipo superan al modelo base, validando el enfoque de modelado individual por equipo.
    \item El modelo BSTS permite capturar cambios en el efecto de localía a lo largo del tiempo, incluyendo el impacto de partidos sin público durante la pandemia.
\end{itemize}

Trabajo futuro incluye:
\begin{itemize}
    \item Incorporación de variables adicionales (fatiga, lesiones, distancia de viaje del visitante)
    \item Modelos de aprendizaje automático (XGBoost, redes neuronales) para mejorar la predicción
    \item Extensión a otras ligas europeas para comparar el efecto de localía entre contextos culturales y estadísticos distintos
\end{itemize}

% -------------------- 8. Referencias --------------------
\begin{thebibliography}{9}

\bibitem{dixoncoles}
Dixon, M. J., \& Coles, S. G. (1997).
\textit{Modelling Association Football Scores and Inefficiencies in the Football Betting Market}.
Journal of the Royal Statistical Society: Series C, 46(2), 265--280.

\bibitem{maher}
Maher, M. J. (1982).
\textit{Modelling Association Football Scores}.
Statistica Neerlandica, 36(3), 109--118.

\bibitem{pollard}
Pollard, R. (2008).
\textit{Home Advantage in Football: A Current Review of an Unsolved Puzzle}.
The Open Sports Sciences Journal, 1(1), 12--14.

\bibitem{karlis}
Karlis, D., \& Ntzoufras, I. (2003).
\textit{Analysis of Sports Data by Using Bivariate Poisson Models}.
Journal of the Royal Statistical Society: Series D, 52(3), 381--393.

\bibitem{scott}
Scott, S. L., \& Varian, H. R. (2014).
\textit{Predicting the Present with Bayesian Structural Time Series}.
International Journal of Mathematical Modelling and Numerical Optimisation, 5(1--2), 4--23.

\bibitem{kaggle}
Panagopoulos, M. (2024).
\textit{English Premier League \& Championship Full Dataset}.
Kaggle. \url{https://www.kaggle.com/datasets/panaaaaa/english-premier-league-and-championship-full-dataset}

\end{thebibliography}

% -------------------- 9. Anexo --------------------
\appendix
\section{Anexo: Código R}

Este anexo incluye el código fuente en R utilizado para la limpieza de datos, estimación de modelos y generación de gráficas. El script principal carga el dataset de Kaggle, ajusta los modelos de Poisson, logístico multinomial y BSTS, y produce las tablas y figuras incluidas en la sección de Resultados.

\subsection*{Estructura del script}

\begin{enumerate}
    \item \texttt{01\_carga\_datos.R}: Lectura y limpieza del CSV.
    \item \texttt{02\_descriptivo.R}: Estadísticas descriptivas y visualización.
    \item \texttt{03\_poisson.R}: Ajuste del modelo de Poisson con efectos por equipo.
    \item \texttt{04\_logit\_multinomial.R}: Modelo logístico multinomial para resultado.
    \item \texttt{05\_bsts.R}: Modelo BSTS para evolución temporal de la ventaja de local.
    \item \texttt{06\_comparacion.R}: Evaluación y comparación de modelos.
\end{enumerate}

\end{document}
